\begin{abstract}
Counterfactual analysis asks a fundamental question in data-driven decision making: what minimal and feasible changes to an input would result in a different outcome? Existing works are largely model-centric, generating synthetic modifications by probing predictive models, while many real-world applications require actionable and trustworthy answers from data. We articulate a vision of \emph{data-centric counterfactual analysis}, where counterfactual reasoning is elevated to a data management primitive and formulated as a constrained, cost-aware and set-oriented retrieval problem. To realize this vision, we introduce \emph{counterfactual querying}, which retrieves feasible, outcome-changing alternatives from observed data under explicit constraints, cost models, and set-level objectives. We develop a two-phase framework that integrates candidate generation that is based on constraints with set-level optimization, guided by the principles of interpretability, cost faithfulness, efficiency, and collective validity. We further present a prototype system, \textbf{DcCA}, demonstrating that counterfactual querying can be implemented efficiently across tabular, temporal, graph, text, and image data. Our results show that counterfactual querying goes beyond similarity search and enables structured, actionable reasoning over data.
\end{abstract}

